\documentclass[11pt, a4paper]{article}

\usepackage[a4paper, margin=2.4cm]{geometry}
\usepackage[utf8]{inputenc}
\usepackage[T1]{fontenc}
\usepackage{amsmath, amssymb, bm}
\usepackage{booktabs, array, tabularx}
\usepackage{enumitem}
\usepackage[colorlinks=true, linkcolor=blue!50!black, urlcolor=blue!50!black]{hyperref}

\setlist{itemsep=0.2em, topsep=0.3em}
\setlength{\parindent}{0pt}
\setlength{\parskip}{0.5em}

\newcommand{\R}{\mathbf{R}}
\newcommand{\X}{\mathbf{X}}
\newcommand{\bt}{\mathbf{t}}
\newcommand{\bc}{\mathbf{c}}

\title{3D-Aware Panorama Stitching with VGGT-Omega}
\author{}
\date{}

\begin{document}
\maketitle

\section{Idea}

Classical stitching estimates 2D homographies between photos and warps them onto a common surface.
This is exact only if the camera rotates about its optical centre or the scene is planar. A camera that
moves, or a scene with depth, produces parallax: objects at different distances shift by different
amounts, and a homography cannot align them, so seams double or cut objects.

Here the photos are not warped to each other. Instead:
\begin{enumerate}
  \item VGGT-Omega gives every photo a camera pose, a depth map and a confidence map, so every
        pixel becomes a point in one shared 3D space.
  \item A virtual camera looks at this 3D scene, and the panorama is its image.
  \item For every panorama pixel, the nearest 3D surface is found with a z-buffer, and the colour is
        read back from the original photos that actually see that surface.
\end{enumerate}
Parallax is resolved by the z-buffer (the nearest surface wins), not by an image warp. The colour always
comes from the original photos, so the output is as sharp as the photos, not as sharp as the model input.

\begin{verbatim}
photos --> VGGT-Omega --> pose, depth, confidence, sky mask
       --> 3D points --> virtual camera + canvas layout
       --> z-buffer (nearest surface per canvas pixel)
       --> backward lookup in every photo + visibility test
       --> exposure gains --> blend (multiband) --> gap fill --> RGBA panorama
\end{verbatim}

\section{Reconstruction}

The $N$ photos are resized to a model resolution (long side 512~px, multiple of 16) and passed to VGGT-Omega
in one batch. For photo $i$ the model returns:
\begin{itemize}
  \item extrinsics $(\R_i, \bt_i)$ (camera-from-world) and intrinsics $\mathbf{K}_i$ from the pose encoding,
  \item a depth map $d_i(u,v)$ (z-depth along the optical axis) and a confidence map $c_i(u,v) \ge 1$.
\end{itemize}
Further per-pixel masks:
\begin{itemize}
  \item \textbf{Depth edges.} A pixel is dropped if the depth range in its $3\times3$ neighbourhood exceeds 5\,\%
        of its depth. These ``flying pixels'' lie between a foreground and a background surface and
        belong to neither.
  \item \textbf{Sky.} An ONNX sky segmenter marks sky pixels. Sky has no valid depth.
\end{itemize}
Depth, confidence and masks are resampled to the colour resolution (long side 1600~px, taken from the
original photo), and $\mathbf{K}_i$ is scaled accordingly. All photos must share one size and orientation.

A pixel $(u,v)$ with depth $d$ is lifted to a world point
\begin{equation}
  \X = \R_i^{\top}\!\left( d\,\mathbf{K}_i^{-1}[u,\,v,\,1]^{\top} - \bt_i \right),
  \qquad \bc_i = -\R_i^{\top}\bt_i .
\end{equation}
Sky pixels have no depth, so they are placed at a fixed very large distance $D_\infty = 10^3 \cdot d_{\mathrm{med}}$
($d_{\mathrm{med}}$ is the median scene depth) along their viewing ray. They then behave like points at infinity:
a rotation-only panorama, with no parallax.

Every other pixel takes part only if its confidence reaches $\tau_i = \max(q_{15}(c_i),\,1)$, where $q_{15}$ is the 15th
percentile of that photo's confidence, and it is not a depth edge.

\section{Virtual camera and canvas}

The canvas is the image of a virtual camera. Two kinds are supported. The choice depends on how the
photographer moved.

\subsection{Perspective canvas (camera that turns)}

A pinhole camera is placed at the mean of the camera centres $\bc_i$. Its frame $\R_p$ has $y$ pointing down
(opposite to the mean camera up vector) and $z$ along the mean horizontal viewing direction of the points
(a circular mean of their longitudes). For a world point $\X$:
\begin{equation}
  [x_p, y_p, z_p]^{\top} = \R_p (\X - \bc),
  \qquad a = \frac{x_p}{z_p}, \quad b = -\frac{y_p}{z_p},
  \qquad u = (a - a_{\min})\,s, \quad v = (b_{\max} - b)\,s .
\end{equation}
This is a standard perspective image, so every straight line in the scene stays straight. The scale $s$ is the
mean focal length in pixels (times a user factor), which keeps the detail of the photos at the canvas centre.
Directions with $z_p/\|\X - \bc\| \le 0.1$ (more than about $84^\circ$ off the axis) are off the canvas.
The extent $a_{\min}, a_{\max}, b_{\min}, b_{\max}$ is the 0.1--99.9 percentile range of all point coordinates plus 1\,\%,
so the canvas is cropped to what the photos cover. The canvas is limited to 50 megapixels.

Limitation: pixels near the border are stretched by $1/\cos^2\theta$ ($\theta$ = angle to the axis), so this mode suits
sets that cover less than about $120^\circ$ horizontally.

\subsection{Plane canvas (camera that walks along a plane)}

For shelves, walls and facades the camera translates, and one pinhole viewpoint would foreshorten the plane
and stretch its far end. Instead the dominant plane is fitted and viewed straight on with an \emph{orthographic} camera.

\textbf{Plane fit.} The points of all photos (sky excluded) go into a RANSAC search (400 hypotheses from 3 random points of a
30\,000-point sample) with inlier tolerance $0.03\,d_{\mathrm{med}}$. A candidate whose normal is more than
$60^\circ$ from the mean viewing direction is rejected, so a floor or ceiling cannot beat the wall in front. The best plane is refined
three times on its inliers: the centroid is the new plane point, and the normal is the eigenvector of the
$3\times3$ covariance with the smallest eigenvalue. If fewer than 15\,\% of the points are inliers, the mode stops
with an error, since there is no dominant plane.

\textbf{Frame.} The normal $\mathbf{n}$ is oriented towards the cameras, and the viewing direction is $\mathbf{f} = -\mathbf{n}$.
The up axis $\mathbf{u}$ is the mean camera up vector projected into the plane, and $\mathbf{r} = \mathbf{f}\times\mathbf{u}$.
With the plane point $\mathbf{o}$:
\begin{equation}
  a = (\X-\mathbf{o})\cdot\mathbf{r}, \qquad b = (\X-\mathbf{o})\cdot\mathbf{u}, \qquad \delta = (\X-\mathbf{o})\cdot\mathbf{f},
\end{equation}
and $u, v$ follow as above with a constant scale $s$ (pixels per world unit). Here $\delta$ is the depth used by the
z-buffer. Since $s$ is constant, one pixel covers the same patch of the plane everywhere and nothing stretches towards the edges.
$s$ is chosen so that one source pixel on the plane covers one canvas pixel: $s = 1/\mathrm{median}(d/f_{\mathrm{px}})$ over the inliers.
Objects in front of the plane still land on the right layer because of the z-buffer.

\section{Z-buffer}

Every usable source pixel of every photo is projected onto the canvas, which gives a canvas position and a depth
(distance to the virtual camera in perspective, $\delta$ in plane mode). Each canvas pixel keeps the \emph{minimum} depth
that reaches it, so the nearest surface wins whichever photo saw it.

Because a source pixel can cover more than one canvas pixel, it is drawn as a square of radius
$\rho = \mathrm{clip}\!\left(\lceil (1.5\,s_{\mathrm{px}} - 1)/2 \rceil,\,0,\,4\right)$, where $s_{\mathrm{px}}$ is the number of canvas pixels one
source pixel covers: $s\,(d/f_{\mathrm{px}})/(D\,z_p^2)$ in perspective ($D$ = distance to the virtual camera) and $s\,(d/f_{\mathrm{px}})$ in plane mode.
Remaining one- or two-pixel holes are closed by two $5\times5$ min-dilations of the depth map.

Larger empty areas (a region that was masked out in every photo) keep no depth. For the lookup in the next step only, their depth is filled
with the average log-depth of valid neighbours (box filters with radii 4, 8, \dots, 256). Such pixels count as ``no strict match''
below.

\section{Colour lookup and visibility}

For each canvas pixel with depth $z$, the 3D point $\X$ is reconstructed and projected into every photo $i$:
$\mathbf{x}_i = \mathbf{K}_i(\R_i \X + \bt_i)$ with camera depth $\hat z_i$. The colour is sampled \emph{bilinearly from the original photo}
(backward warping), and depth, confidence, sky and edge are sampled nearest. Photo $i$ is a \emph{strict} observer of the pixel if
\begin{equation}
  \mathbf{x}_i \text{ lies inside the photo}, \quad |d_i(\mathbf{x}_i) - \hat z_i| \le 0.08\,\hat z_i, \quad c_i(\mathbf{x}_i) \ge \tau_i, \quad \text{not sky, not edge},
\end{equation}
that is, the photo sees \emph{the same surface} that the z-buffer chose. A photo that sees something nearer
(an occluder) or farther is rejected for this pixel. For sky pixels (depth $\ge D_\infty/2$) the test is just ``the photo marks this pixel as sky''.

Each strict observer gets the weight
\begin{equation}
  w_i = \bigl(\mathrm{clip}\bigl(\tfrac{c_i - 1}{c^{90}_i - 1},\,0,\,1\bigr) + 0.05\bigr)\cdot \mathrm{smoothstep}\!\bigl(\tfrac{\text{border distance}}{0.2\cdot\min(H,W)}\bigr),
\end{equation}
where $c^{90}_i$ is the 90th percentile of the photo's confidence. Confident pixels count more, and weights fall to zero smoothly
at the photo border so that seams do not sit on image edges. Sky pixels get a fixed confidence of 0.5.

\section{Exposure matching}

Photos taken a moment apart can differ in exposure and white balance. Per photo and colour channel a gain $g_i$ is
estimated on a low-resolution probe render (at most 1200~px wide), over the pixels that at least one other photo also observes:
\begin{equation}
  g_i = \frac{\sum_p w_{ip}\,\bar I_{-i}(p)}{\sum_p w_{ip}\,I_i(p)},
  \qquad \bar I_{-i}(p) = \frac{\sum_{j\ne i} w_{jp}\,g_j I_j(p)}{\sum_{j\ne i} w_{jp}} .
\end{equation}
Three rounds are run, the gains are divided by their mean over photos, and clipped to $[0.5, 2]$. The gains are applied at
render time, so the original pixels are not altered.

\section{Blending}

\textbf{Multiband mode (default).} Averaging all observers blurs fine detail when two photos are slightly misaligned, because two copies of the same
edge are mixed. Multiband blending avoids this by choosing one photo per pixel and hiding the seam in a pyramid:
\begin{enumerate}
  \item Each weight map is smoothed with a Gaussian of $\sigma = \max(3,\,0.01\min(H,W))$. Each pixel is assigned to the strict
        observer with the highest smoothed weight, which gives a hard label map $M_i$.
  \item Every photo (after its gain) gets a Laplacian pyramid $L_k(I_i)$ and its mask a Gaussian pyramid $G_k(M_i)$, with
        $K = \min(6,\,\max(1,\lfloor\log_2(\min(H,W)/24)\rfloor))$ levels. Where a photo has no valid content it is painted
        with the plain weighted average, so no black bleeds into the bands.
  \item Each level is combined as $B_k = \sum_i L_k(I_i)\,G_k(M_i) \big/ \sum_i G_k(M_i)$ and the pyramid is collapsed.
\end{enumerate}
Low frequencies therefore blend over a wide band (no visible brightness step), and fine detail blends over a narrow band, so each edge comes from a single photo.

\textbf{Average mode.} The weighted average $\sum_i w_i g_i I_i / \sum_i w_i$. It is used instead of multiband if photos
$\times$ canvas pixels exceeds $1.5\cdot 10^8$.

\section{Gap filling}

If no photo passes the strict test at a pixel (all observers are masked by confidence, edge or depth disagreement), it would stay empty
although some photo does see that direction. For these pixels only, a fallback colour is used: any photo whose image contains the
pixel contributes with its border weight, \emph{without} the depth, confidence and edge tests. The fallback is accumulated separately
from the strict result, so it never dilutes a pixel that has a strict match. Such pixels are reported as \texttt{filled\_from\_fallback}.

What is still empty (outside every photo, such as the corners of a perspective canvas) is left transparent. Gaps up to a few pixels wide are closed by a
$5\times5$ morphological closing and Telea inpainting. The result is saved as an RGBA PNG whose alpha is 0 where no photo covers the canvas.

\section{Parameters}

\begin{center}
\small
\begin{tabularx}{\textwidth}{@{}l l X@{}}
\toprule
\textbf{Parameter} & \textbf{Default} & \textbf{Effect} \\
\midrule
projection & perspective & \texttt{perspective} for a camera that turns, \texttt{plane} for one that walks along a wall or shelf \\
scale & 1.0 & Canvas pixels per source pixel at the canvas centre (1.0 keeps the photo detail) \\
blend & multiband & \texttt{multiband} or \texttt{average} \\
fill\_gaps & on & Fallback colour where no photo passes the strict tests \\
conf\_percentile & 15 & Share of the least confident pixels of each photo that is dropped \\
depth\_tolerance & 0.08 & Relative depth disagreement that still means ``same surface'' \\
feather & 0.2 & Border ramp, as a fraction of the shorter photo side \\
edge\_rtol & 0.05 & Depth range in a $3\times3$ window (relative) above which a pixel is a flying pixel \\
exposure & on & Per-photo, per-channel gain \\
use\_sky & on & Place sky at a very large distance \\
max\_megapixels & 50 & Canvas size cap \\
model resolution & 512 & Long side of the network input \\
colour resolution & 1600 & Long side of the photos used for colour \\
\bottomrule
\end{tabularx}
\end{center}

\section{Limitations}

\begin{itemize}
  \item \textbf{Accuracy is bounded by the model.} Pose and depth come from a $512$~px network pass. A pose error of a fraction of a degree
        or a depth error on a thin object shows up as a small misalignment or ghosting. Nothing in the current pipeline
        corrects the model output with the photos.
  \item \textbf{Fallback pixels are less reliable.} They skip the depth test, so near objects seen from different positions can ghost there.
        Switch \texttt{fill\_gaps} off to see how many pixels depend on it.
  \item \textbf{Perspective mode} stretches the borders beyond about $120^\circ$ of horizontal view and fails beyond $170^\circ$. Use plane mode for a walking camera.
  \item \textbf{Plane mode} assumes one dominant plane that faces the cameras. Corners, corridors and scenes without a plane are not supported and give an error.
  \item \textbf{No lens model.} The pinhole camera of the network is used as is. Visible lens distortion is not corrected.
  \item \textbf{Moving objects} appear in one photo only or in several positions. The pipeline does not detect motion.
  \item \textbf{All photos must have the same size and orientation.}
\end{itemize}

\section{Code}

\begin{itemize}
  \item \texttt{panorama/stitch.py}: reconstruction, layouts, z-buffer, lookup, exposure, blending, gap filling.
  \item \texttt{panorama/app.py}: Gradio interface. The model output of a photo set is cached, so changing render options re-renders in seconds without running the network.
  \item \texttt{panorama/test\_panorama.py}: synthetic checks on a textured sphere and a textured wall.
\end{itemize}

\end{document}
